\documentclass[11pt]{article}
\usepackage[margin=25mm]{geometry}
\usepackage{fontspec}
\setmainfont{Latin Modern Roman}
\usepackage{amsmath,amssymb,mathtools}
\usepackage{unicode-math}
\setmathfont{Latin Modern Math}
\usepackage{xurl,hyperref}
\hypersetup{colorlinks=true,urlcolor=blue}
\setcounter{secnumdepth}{0}
\setlength{\parindent}{0pt}
\setlength{\parskip}{5pt}
\setlength{\emergencystretch}{3em}
\newcommand{\sref}[2]{\hyperlink{source:#1:#2}{[S#2]}}
\newcommand{\cataloguescope}[1]{\paragraph{Recorded target.}#1}
\newcommand{\R}{\mathbb R}
\newcommand{\C}{\mathbb C}
\newcommand{\N}{\mathbb N}
\newcommand{\Z}{\mathbb Z}
\newcommand{\Q}{\mathbb Q}
\newcommand{\F}{\mathbb F}
\newcommand{\eps}{\varepsilon}
\begin{document}
\section{One-Dimensional Fermi Gas with Attractive Interaction}\label{statement}

\subsection{Definitions and mathematical statement}

\paragraph{Shared notation.}$\mathbb N=\{1,2,\ldots\}$; the standard number fields have their usual meanings. All additional conventions are stated locally.


\paragraph{Definitions and notation.}Let $\psi_s(x)$ be fermionic fields with spin index $s$, and let the one-particle density matrix be $\rho_{ss'}(x,y)=\langle\psi_s^\dagger(x)\psi_{s'}(y)\rangle$. For a translation-invariant state it depends on $x-y$. The source considers a weakly interacting one-dimensional gas, first in an interval with periodic boundary conditions, then in the thermodynamic limit and the zero-temperature limit. Mass, chemical potential, pair potential and coupling determine the model. The target concerns spinful attractive interactions, including how decay exponents depend on the coupling.

\paragraph{Catalogue wording retained for traceability.}
Determine the large-distance behavior of the one-particle reduced density matrix for a one-dimensional Fermi gas with spin and attractive interaction.
\paragraph{Scope and source qualification.}The catalogue target is the one-dimensional spinful gas. The source additionally discusses strip and two-row lattice variants as related first steps, without identifying them with this canonical record. Smooth finite-range potentials are mentioned as a context for known perturbative results, not imposed as a replacement for the original target.
\subsection{Short English statement}
Determine the long-distance one-particle correlations of a weakly attractive one-dimensional spinful Fermi gas.
\subsection{Research attempt}
\paragraph{Model, date and investigated scope.}
OpenAI GPT-6 Astra (configured model \texttt{gpt-6-astra}), ultra reasoning; generated 11 October 2026. Computed the exact free spinful kernel and demonstrated the nonuniformity of exponent perturbation at large distance. Attractive correlations, exponents and possible nonperturbative scales remain unresolved.

\paragraph{Sources and literature check.}
Read the supplied complete definition, catalogue wording, scope qualifications and source list for this canonical ID. The original-source citations are inherited as supplied references; no fresh primary-source or novelty search was performed for this attempt. No current open-status claim is inferred from this calculation.

\paragraph{The zero-coupling benchmark.}For a translation-invariant free gas with two spin states at zero temperature, suppose each spin occupies momenta \([-k_F,k_F]\), \(k_F>0\). In the Fourier convention \(\psi(x)=(2\pi)^{-1/2}\int e^{ikx}a(k)\,dk\), the spin-diagonal one-particle kernel is
\[\rho_{ss'}(0,r)=\delta_{ss'}\frac1{2\pi}\int_{-k_F}^{k_F}e^{ikr}\,dk
=\delta_{ss'}\frac{\sin(k_Fr)}{\pi r},\]
with the continuous value \(k_F/\pi\) at \(r=0\). This follows immediately from the occupied-mode projection; off-diagonal spin components vanish for the chosen unpolarized state. It fixes both the oscillation and the \(r^{-1}\) envelope at the noninteracting endpoint.
\paragraph{Why finite-order perturbation is insufficient.}A candidate coupling-dependent algebraic envelope might have the form \(r^{-1-\alpha(g)}\), with attractive coupling \(g<0\) and \(\alpha(g)=ag^2+O(g^3)\). Its formal expansion begins
\[r^{-1-\alpha(g)}=r^{-1}\left(1-ag^2\log r+\cdots\right).
\]
For fixed nonzero \(g\), the purported second-order correction ceases to be small when \(\log r\) is comparable to \(|g|^{-2}\). This is an exact warning about the model expression: expanding \(e^{-\alpha(g)\log r}\) requires \(|\alpha(g)\log r|\ll1\). Likewise an arbitrarily small positive inverse correlation length \(m(g)\) in \(r^{-1}e^{-m(g)r}\) changes the fixed-\(g\) large-distance law, however closely it approximates the free expression on bounded intervals. These are competing ansatz tests, not a claim that either is the actual attractive-gas law.
\paragraph{Gap and limits.}The first missing step is a uniform-in-distance control of interacting correlations after the thermodynamic and zero-temperature limits, including possible spin-sector scales. No such bound is derived here. The spinful attractive target cannot be replaced by a free gas or a spinless power law, and the source's strip variants remain related questions rather than substitutes. The benchmark and nonuniform-expansion calculation establish no new attractive-coupling exponent.

\paragraph{Verification and adversarial audit.}
Checked Fourier normalization, r=0, spin diagonal structure and the g squared log r expansion domain; did not assert any unproved interacting decay ansatz. This was an LLM adversarial audit, not independent human verification; no theorem prover was used. Prompts A then B applied to this ID: Checked Fourier normalization, r=0, spin diagonal structure and the g squared log r expansion domain; did not assert any unproved interacting decay ansatz. Scope remains limited: Control the interacting one-particle correlation uniformly in distance at fixed negative coupling after the required limits.

Coordinator checked Fourier normalization, the r=0 value and the nonuniform logarithmic expansion. No interacting exponent or correlation-length claim is made. This additional check was model-assisted, not independent human verification.

\paragraph{Tools, process and contributions.}
Python 3 and local shell used to read supplied files and serialize this manuscript; mathematics derived symbolically in the text, with no numerical experiment or formal proof assistant. Contribution prompt A was followed by prompt B. GPT-6 Astra with ultra reasoning drafted the research and performed its stated self-audit; any additional coordinator contribution is identified in the verification and audit. A GPT-6-based Codex coordinator assembled the manuscript and checked repository integration; its exact variant and reasoning configuration were not exposed. No human mathematical contribution, independent human proof check or formal proof certification is supplied. Complete derivations and the consulted source references are included here; private reasoning transcripts are not published.

\paragraph{Final status and first remaining gap.}
\textbf{UNRESOLVED}. Control the interacting one-particle correlation uniformly in distance at fixed negative coupling after the required limits.

\section{Completion Estimate}
\noindent\textbf{COMPLETION: 0\%}
The calculation is an elementary benchmark or conditional reduction; it establishes no claimed new portion of the complete source target.

\subsection{Sources}
\begin{itemize}
\item[\hypertarget{source:statement:1}{S1}] ULAM.AI and the UnsolvedMath Contributors, UnsolvedMath, catalogue entry. Catalogue curation under CC BY 4.0; definitions and English exposition added locally.
\url{https://huggingface.co/datasets/ulamai/UnsolvedMath}.
\item[\hypertarget{source:statement:2}{S2}] Giovanni Gallavotti, original contribution to Open Problems in Mathematical Physics (1998), edited by Michael Aizenman.
\url{https://web.math.princeton.edu/~aizenman/OpenProblems_MathPhys/9807.FermiGas.html}.
\end{itemize}
\label{end:statement}
\end{document}
